\documentclass[11pt]{article}
\usepackage[T1]{fontenc}\usepackage{lmodern,amsmath,amssymb,amsthm,booktabs}
\usepackage[margin=1in]{geometry}
\newtheorem{theorem}{Theorem}\newtheorem{lemma}{Lemma}
\title{A rational post-half-mass interval for the weak-produced input sign}
\author{Analytic increment for independent review}\date{Version 1, 2 November 2026}
\begin{document}\maketitle
\begin{abstract}
On the selected orbit the input gap \(z=X-Y\) is positive through
\(s_1=1/1000\) in the weak half-mass clock, uniformly for
\(13/20\le\rho\le7/10\). We give lower bounds for both \(T=-2Y'\)
and \(z\) on the new interval. On the intersection of WIT's learned
window with \(s\le s_1\), the limiting weak-produced input contribution
has a positive uniform minimum and the initialized contribution has
the same sign for small noise. The boundary estimates below leave
the part of the window after \(s_1\) unresolved.
\end{abstract}
\section{Domains and clocks}\label{an01phsv1:domains}
Write \(x,y,u,v\) for the coherent comparator, \(\lambda=\rho^2\),
\(n=(1+e^{-\lambda s})^{-1}\), \(z=x-y\), \(A=\rho K(\rho x)\),
\(P=\Phi(\rho x)\), \(Q=\Phi(u)\), \(T=y+nK(u)-(1-n)A\).
Here \(K(t)=\phi(t)+t\Phi(t)\), \(K'=\Phi\). The exposure in
the producer formula is \(\zeta-x\), distinct from \(A\).
The following are the source domains used throughout.
\begin{itemize}
\item WIT \texttt{orbit} and \texttt{Tequation} concern the selected
coherent leading branch of strong mass one, at finite \(s\) with
\(0<n<1\). The identity for \(T'\) has this domain. The lower bound
\(T\ge (5/62)n(1-n)\) and the region
\(-n<y<a<2/5,\ y<x<3/4,\ 0<u<1,\ 0<v<11/10\) in
\texttt{regiontheorem} require the WIT ratio box and \(s\le0\)
(\(n\le1/2\)). The elementary Gaussian bounds apply throughout that box.
\item P \texttt{limitreversal} gives finite-time continuation and
\(v>0,\ v>x,\ v>y\) on the positive weak-growth branch for
\(0<\rho<1\). Its probe reversal statement additionally assumes
\(\zeta>a_\rho\). Only its probe-independent invariant and growth
estimates are used here.
\item WIT \texttt{branchsmooth} supplies continuous parameter
dependence on fixed finite windows of the selected branch.
\texttt{commonwindows} fixes the common compact \(J\), with
\(s_a=-(9/4)\log3<s_*\) at the central ratio and the window inside the
central pre-hit interval. \texttt{uniformtheorem} supplies the mass,
probe-gap and rate margins only on \(\mathcal P_*\times J\).
Its construction does not put \(\sup J\) before half mass or before
the rational endpoint below.
\item CPD's table and help theorem concern the original aligned balanced
initialization \(S_0=q^{10}\) on \(\mathcal P_*\times J\), for small
positive \(q\) after fixed constants. The rates use the selector
convention (classical at regular states, a.e.\ for the time derivative).
The phase is \(s=\mu_1^2 t_{\rm phys}-\tau_h\). CPD uses
\(R^{(p)}=R_p/2\) and the full reaction \(h_p=2g_p\).
\end{itemize}

\section{Crossing test and the attempted extensions}\label{an01phsv1:cross}
The orbit gives \(2y'=-T\), \(2x'=\lambda P[n(v-y)-z]\).
At any \(z=0\) with \(n\ge1/2\), keeping the signed term gives
\[
 2z'=\lambda n\Phi(\rho y)(v-y)+y+nK(u)-(1-n)\rho K(\rho y).
\]
For \(y\ge0\) group the \(y\) term with the linear part of \(K\):
\(y[1-(1-n)\lambda\Phi(\rho y)]\ge0\). For \(y<0\),
\(K(\rho y)\le K(0)\). In either case
\(\rho\phi(\rho y)\le\rho\phi(0)<(7/10)(2/5)=7/25\).
Dropping the strictly positive first term and using \(1-n\le n\)
therefore proves
\begin{equation}\label{an01phsv1:reviewer}
 2z'>n[K(u)-7/25]-y_-,\qquad y_-=\max(-y,0).
\end{equation}
This includes \(y=0\) and \(n=1/2\). At \(u=0\) the exact formula is
\(2u'=(1-n)\phi(0)-y/2\). Both proposed local tests are valid.
In particular \(u,y\ge0\) suffice at a crossing since \(K(0)>3/8\).

\paragraph{Approach A.} To propagate \(T>0\) and \(z>0\) together,
consider a first \(T=0\) with \(z\ge0\). WIT's exact equation there is
\[
 \frac{T'}{n(1-n)}
 =A(\lambda-Q^2/2)+\frac\lambda2(K(u)+\phi(u))
 +\frac12Q\phi(u)+\frac{n}{2(1-n)}(K(u)-u)Q^2
 -\frac{\lambda P x'}n .
\]
Let \(R=\max(|x|,|y|,|u|,|v|)\). For all \(s\ge0\), P's Dini
bound gives \(R(s)\le R_s=(11/10+s/5)\exp[(1-\lambda)\log2/(2\lambda)]\)
by integrating \(1-n\) from zero to infinity.
Then \(P\le\Phi(\rho R_s), Q\le\Phi(R_s),v-y\le2R_s\).
Even retaining the positive forcing terms, the sufficient inward test
obtained from \(x'\le\lambda nP(v-y)/2\) needs their sum plus
\(A(\lambda-Q^2/2)\) to exceed \(\lambda^2P^2(v-y)/2\).
The half-learning estimates \(Q\le9/10,P\le71/100,v-y\le8/5\)
do not follow from these forward bounds. For example a direct uniform
lower estimate (discarding positive terms) is
\[
 -\tfrac{31}{400}\rho K(\rho R_s)
 +\tfrac{\lambda}{2}K(-R_s)-\lambda^2 R_s ,
\]
which is negative (\(K(-R_s)<2/5,\ R_s>1,\lambda>2/5\)).
This estimate leaves \(T'>0\) unpaid. At a fixed state
\(\partial_n T=K(u)+A>0\), but along the orbit \(A'=\lambda Px'\);
the last term in the exact test is precisely the cost of moving \(x\).
The fixed-state derivative does not settle its sign.

\paragraph{Approach B.} A first-exit proof using the nonnegative quadrant
needs lower entry bounds and inward velocities. WIT supplies only
\(y(0)>-1/2\); its upper bound \(y<a\) does not provide
\(y<2(1-n)\phi(0)\) as \(n\) grows. At \(y=0\) one additionally needs
\((1-n)A\ge nK(u)\) for inward \(y'\). Even \(u=x=y=0\) and
\(n\ge1/2\) give the reverse strict inequality.
For a crossing with \(y<0\), retaining the positive term instead
requires
\begin{equation}\label{an01phsv1:missing}
 \lambda n\Phi(\rho y)(v-y)+y+nK(u)>(1-n)\rho K(\rho y).
\end{equation}
To see the limitation of the pointwise bounds, take a boundary test
state \(x=y=-1/5,\ u=0,\ v=1/20,\ n=9/16\). It satisfies all three
forward \(v\) inequalities, \(u\ge0,\ y>-n\), and even
\(y+nK(u)>-1/5+27/128>0\).
The left side of \eqref{an01phsv1:missing} is strictly less than
\(9/256+1/40\) using \(\lambda<1/2,\Phi(\rho y)<1/2\).
By convexity \(K(t)\ge K(0)+t/2\), so its right side exceeds
\((7/16)(13/20)(3/8-7/100)>9/256+1/40\).
Thus these inequalities alone do not close the crossing test.
This test state is not asserted to occur on the orbit. Neither failed
extension gives evidence of an actual sign change.

\section{A rational interval from the half-mass margin}\label{an01phsv1:local}
\begin{theorem}\label{an01phsv1:theorem}
For the whole ratio box and \(0\le s\le s_1=1/1000\),
\[
 R(s)<2,\qquad T(s)\ge5/248-12s_1>0,\qquad
 z(s)\ge e^{-s/4}\frac{15}{3968}>0 .
\]
Consequently \(W_{\rm in}=(\lambda/2)n z(v-y)P>0\) through \(s_1\),
including the entire preceding finite selected branch.
\end{theorem}
\begin{proof}
On \([-1,0]\), \(1/4<n\le1/2\) since \(e^\lambda<3\).
WIT's half-learning region and the exact \(z\) equation imply
\[
 z'+z/4\ge T/2\ge(5/124)n(1-n)\ge5/992 .
\]
Here \(\lambda P/2<1/4,\ z>0,\ v-y>0\).
Integration over that unit interval gives \(z(0)\ge(3/4)(5/992)\)
using \(e^{-1/4}\ge3/4\). Also \(T(0)\ge5/248\).
These bounds use no learned-window margins.

At zero \(R(0)\le11/10\). On the closed box \(R\le2\),
\(K(u)<12/5,\ A< (7/10)(2/5+7/5)<13/10,\ |n'|<1/8\).
Each coordinate speed is at most 4: the brackets in \(2x',2y',2v'\)
are bounded in magnitude by \(4\lambda,\ 2+12/5+13/10,\ 13/10+2\lambda\),
respectively, and in \(2u'\) by \(2/5+2\lambda+4\).
Thus before any first box exit \(R(s)\le11/10+4s\), strictly
below 2 through \(s_1\). Differentiating the definition of \(T\)
within the reached box gives
\[
 |T'|\le |y'|+|n'|(K(u)+A)+nQ|u'|+(1-n)\lambda P|x'|
 <4+(12/5+13/10)/8+4+2<12 .
\]
It follows that \(T\) is strictly positive through the closed endpoint.
Finally
\[
 2z'+\lambda Pz=\lambda nP(v-y)+T>0 .
\]
The integrating factor gives the stated \(z\) bound. For earlier times
apply WIT's region and the forward \(v-y\) invariant directly.
\end{proof}

\section{Original-flow conclusion and open interval}\label{an01phsv1:transfer}
Put \(J_1=J\cap(-\infty,s_1]\). If this set is nonempty it is compact.
WIT branchsmooth and Theorem \ref{an01phsv1:theorem} give
\(m_1=\min_{K_\rho\times J_1}W_{\rm in}>0\).
Keep CPD's constants first, fix \(s_1\) and this minimum, then take
q small enough that CPD's uniform producer error \(Cq^{1/20}<m_1/2\).
The exact selected-weak-produced core input contribution to source 2
is then at least \(q^2m_1/2\) on \(\mathcal P_*\times J_1\).
An empty \(J_1\) entails no transfer claim.

For the initialized \(S_0=q^{10}\) population, CPD gives on all of J
\[
 q^{-2}D_1^\tau=-\tfrac12\mathcal A_q\mathcal L_q+O(q^{1/20}),\qquad
 q^{-2}D_2^\tau=T_{\rm in}+W_{\rm in}+(1-n)T_{\rm out}
                   +O(q^{1/20}).
\]
Here the exact lag is positive for small q by CPD help,
\(T_{\rm in}=-\lambda z^2P/2\),
\(T_{\rm out}=\rho(\zeta-x)K(\rho x)/2>0\), and
\(W_{\rm out}=-nT_{\rm out}<0\).
CPD retains all fixed-label producers, radial and non-core charges.
On \(J_1\) the weak-produced input term has the positive sign proved
above. On any part of J where \(z\le0\), both limiting input terms
are nonpositive, so the positive limiting \(d_2\) is carried by
\((1-n)T_{\rm out}\ge d_2\ge5p_*/8\). This implication does not
assert existence of such a subwindow.

The scope of the updated mechanism is therefore target/core and
weak-produced input rotation together with positive net output rotation
of the strong carrier. The weak input sign on \(J\cap(s_1,\infty)\)
still needs \eqref{an01phsv1:missing} at any first zero, or a
window-placement bound. Physical rates multiply the normalized rates
by \(\mu_1^2\) in the matched clock \(s=\mu_1^2t_{\rm phys}-\tau_h\).
The canonical and bridge regimes need their own initialized state
and producer errors below the actual signed lag margin.
\end{document}
